

\centerline{{\bf Abstract}}
\bigskip \noindent


%=======================Begin of the abstract ===========================
%% Explain the title

A Laurent-Ore algebra~$\LO$ over a field~$F$ is a mathematical
abstraction of common properties of linear partial differential
and difference operators. A linear (partial) functional system is
of the form~$A(\vz)=0$ where~$A$ is a matrix over~$\LO$ and~$\vz$
is a vector of unknowns. Typically, it is a system consisting of
linear partial differential, shift and~$q$-shift operators, or any
mixture thereof.


%% Link modules and systems

We associate a linear functional system~$A(\vz)=0$ to
an~$\LO$-module~$M_A$, which is called the module of formal
solutions. For our purpose, the %linear
dimension of an~$\LO$-module is  defined to be the dimension of
the module as a vector space over~$F$. A system~$A(\vz)=0$ is said
to be~$\pa$-finite if~$M_A$ has finite dimension.


A Picard-Vessiot extension for a~$\pa$-finite system~$A(\vz)=0$ is
a ring containing ``all'' solutions of~$A(\vz)=0$. We prove the
existence of Picard-Vessiot extensions for all~$\pa$-finite linear
functional systems and show that the dimension of the solution
space of a~$\pa$-finite system equals the dimension of its module
of formal solutions.

The Gr\"obner basis techniques for left ideals in Ore algebras are
extended to left submodules over Laurent-Ore algebras.  This
extension enables us to determine whether a linear functional
system is~$\pa$-finite.
% compute the linear dimension of a module
%of formal solutions.

We present an algorithm for finding all submodules of
an~$\LO$-module with finite dimension. This algorithm allows us to
find all ``subsystems'' whose solution spaces are contained in
that of a given~$\pa$-finite linear functional system.


%===============  End of the abstract ===========================
